% ══════════════════════════════════════════════════════════════════════
% § 4   Boundary analysis
% ══════════════════════════════════════════════════════════════════════

\subsection{The boundary one-loop coefficient
  \texorpdfstring{$b_{0}^{(b)}(g_{a}^{\star})$}{b0(b)(ga*)}}
\label{sec:boundary-b0}

Recall from Section~\ref{sec:setup-boundary} that the stability of
the axis-$a$ SCFT at
$(g_{1},g_{2})=(\delta_{a1}\,g_{a}^{\star},\delta_{a2}\,g_{a}^{\star})$
against turning on the orthogonal coupling is determined by the
effective one-loop coefficient of the decoupled node $b\ne a$:
\begin{equation}
  b_{0}^{(b)}(g_{a}^{\star})\;=\;T(\adj_{b})\;+\;\sum_{i\text{ at node }b}
    T_{G_{b}}(r_{i})\,\bigl(R_{i}^{(a\star)}-1\bigr)\,.
  \label{eq:b0-boundary-rep}
\end{equation}
At leading $N$ this has the linear form
\begin{equation}
  b_{0}^{(b)}(g_{a}^{\star})\;=\;\alpha_{b}^{(a\star)}\,N
  +\gamma_{b}^{(a\star)}\;-\;\tfrac{1}{2}\,\Nf^{(b)}\,,
  \label{eq:b0-linear}
\end{equation}
where the $-\Nf^{(b)}/2$ term arises because the fundamentals on
the decoupled node $b$ retain their free R-charge $R=2/3$ and
contribute $T_{G_{b}}(\fund)\cdot(2/3-1)=-1/4$ per pair to the
leading-$N$ sum --- precisely the UV value.  Absorbing the Veneziano
shift, the condition ``axis $a$ is marginally IR-unstable'' reads
$b_{0}^{(b)}(g_{a}^{\star})=0$, i.e.
\begin{equation}
  \Nf^{(b)}\;\le\;\alpha_{b}^{(a\star)}\,N
  +2\gamma_{b}^{(a\star)}\,,
  \label{eq:boundary-inequality}
\end{equation}
in direct analogy with the UV AF bound
$\Nf^{(b)}\le 2\alpha^{(b)}N+2\gamma^{(b)}$ obtained by setting the
bifundamental anomalous dimension to zero.

The leading-$N$ coefficient $\alpha_{b}^{(a\star)}$ is the central
datum of boundary analysis: its \emph{sign} determines, in the
non-Veneziano limit $\Nf^{(b)}=0$, whether axis $a$ is IR-unstable:
\begin{equation}
  \boxed{\;
  \alpha_{b}^{(a\star)}>0\;\Longleftrightarrow\;b_{0}^{(b)}(g_{a}^{\star})>0
  \;\Longleftrightarrow\;\text{axis $a$ is IR-unstable}
  \;\Longleftrightarrow\;\text{flow released into interior.}
  }
\end{equation}

\subsection{The nine equal-rank classes}
\label{sec:boundary-table}

Table~\ref{tab:b0-boundary} collects the UV and boundary one-loop
bounds for the eight non-Veneziano equal-rank classes with a
well-defined leading-$N$ $a$-maximum and non-empty matter at each
node.  Class~877 and the 20 flat-direction theories belong to the
nonSCFT umbrella of Section~\ref{sec:class-special} and are omitted;
for class~877 specifically the empty-node morphology places the
entire matter sector on bifundamentals, so the sub-quiver decoupling
limit required for $\alpha_{b}^{(a\star)}$ needs a separate treatment.
Classes~$1$ and~$9$, although also in the nonSCFT umbrella, admit
well-defined formal leading-$N$ $\alpha_{b}^{(a\star)}$ and are
retained in the table with this caveat.  The
$\alpha$-coefficients are class-invariants (independent of the choice
of representative within a class, as we have verified for each
class in Table~\ref{tab:class-reps}).

\begin{table}[h]
  \centering
  \small
  \renewcommand{\arraystretch}{1.18}
  \setlength{\tabcolsep}{4pt}
  \begin{tabular}{clllll}
    \toprule
    Class & UV bound node 1
          & UV bound node 2
          & Boundary-1 bound on $\Nf^{(2)}$
          & Boundary-2 bound on $\Nf^{(1)}$
          & Phase \\
    \midrule
    $1$   & $\Nf^{(1)}\!\le 2N{-}1$    & $\Nf^{(2)}\!\le 2N{-}1$
          & $\Nf^{(2)}\!\le N{-}1$     & $\Nf^{(1)}\!\le N{-}1$
          & Fig.~5$^{\ast}$ \\
    $2$   & $\Nf^{(1)}\!\le 2N{-}1$    & $\Nf^{(2)}\!\le N{-}5$
          & $\Nf^{(2)}\!\le 0.2188\,N{-}5$ & $\Nf^{(1)}\!\le 1.2188\,N{-}1$
          & Fig.~5 \\
    $4$   & $\Nf^{(1)}\!\le 3$         & $\Nf^{(2)}\!\le 2N{-}1$
          & $\Nf^{(2)}\!\le 1.2431\,N{-}1$ & $\Nf^{(1)}\!\le -0.7569\,N{+}3$
          & mixed \\
    $9$   & $\Nf^{(1)}\!\le \tfrac34 N{+}c_{1}$ & $\Nf^{(2)}\!\le \tfrac34 N{+}c_{2}$
          & $\Nf^{(2)}\!\le \tfrac34 N{+}c_{2}$ & $\Nf^{(1)}\!\le \tfrac34 N{+}c_{1}$
          & Fig.~5$^{\ast}$ \\
    $16$  & $\Nf^{(1)}\!\le 1$         & $\Nf^{(2)}\!\le N{-}1$
          & $\Nf^{(2)}\!\le 0.4270\,N{-}1$ & $\Nf^{(1)}\!\le -0.5730\,N{+}1$
          & mixed \\
    $21$  & $\Nf^{(1)}\!\le N{+}2$     & $\Nf^{(2)}=0$ (marg.)
          & $\Nf^{(2)}\!\le -0.3246\,N$    & $\Nf^{(1)}\!\le 0.6754\,N{+}2$
          & mixed$^{\dagger}$ \\
    $23$  & $\Nf^{(1)}\!\le N{+}1$     & $\Nf^{(2)}\!\le 3$
          & $\Nf^{(2)}\!\le -0.1468\,N{+}3$ & $\Nf^{(1)}\!\le 0.8532\,N{+}1$
          & mixed \\
    $44$  & $\Nf^{(1)}\!\le 3$         & $\Nf^{(2)}\!\le 3$
          & $\Nf^{(2)}\!\le 3$         & $\Nf^{(1)}\!\le 3$
          & marginal \\
    \bottomrule
  \end{tabular}
  \caption{UV and boundary one-loop bounds for the eight equal-rank
    non-Veneziano classes with defined leading-$N$ $a$-maximum
    (class~877 omitted; see Section~\ref{sec:class-special}).
    ``Phase'' anticipates the classification of
    Section~\ref{sec:phases}: \emph{Fig.~5} means
    $\alpha_{1}^{(2\star)},\alpha_{2}^{(1\star)}>0$ (both axes
    IR-unstable, interior SCFT attracts); \emph{mixed} means one
    $\alpha$ positive, one negative.  $^{\ast}$Classes 1 and 9
    belong to the nonSCFT umbrella of
    Section~\ref{sec:class-special}: they are formally in the Fig.~5
    sign pattern at leading $N$ but their interior R-charges are below
    the unitarity bound; the Fig.~5 phase assignment for both should
    be understood as reflecting the leading-$N$ sign pattern only.  For class 9 the UV and boundary bounds coincide
    ($\alpha_{b}^{(a\star)}=\alpha^{(b)}=3/4$) because every matter
    field sits at $R=1/2$ — the $R$-independent and
    $R$-shift contributions to $\alpha$ cancel by design.  The
    intercepts $c_{1},c_{2}$ depend on the representative; all
    representatives share $\alpha=3/4$.  $^{\dagger}$Class 21 has
    $\alpha^{(2)}=0$ at the UV — node 2 is marginal at leading $N$;
    the boundary $\alpha$ on axis 2 is negative and well-defined.
    $^{44}$Class 44 is leading-$N$ marginal at both boundaries:
    $\alpha$ vanishes on both sides.}
  \label{tab:b0-boundary}
\end{table}

\paragraph{Reading the table.}
At $\Nf^{(a)}=0$ (non-Veneziano limit), the sign of the boundary
coefficient $\alpha_{b}^{(a\star)}$ alone determines whether axis $a$
is IR-unstable ($>0$) or IR-stable ($<0$) in the large-$N$ limit.
The values $\alpha_{b}^{(a\star)}$ are obtained from the symbolic
$a$-maximisation of each class by substituting the bifundamental
R-charge $R_{\text{bif}}^{(a\star)}$ into the general formula
\eqref{eq:b0-boundary-rep}; for equal-rank classes the result is
equivalently expressible in the form \eqref{eq:boundary-inequality}
with $\alpha,\gamma$ rational functions of $R_{\text{bif}}^{(a\star)}$.

\subsection{RG-topology classification of the nine classes}
\label{sec:boundary-phases}

Each pair $(\alpha_{1}^{(2\star)},\alpha_{2}^{(1\star)})$ of
boundary coefficients determines the RG topology of the class at
$\Nf=0$.  The possible sign patterns, with names following
\cite{IW2005}, are:
\begin{center}
  \begin{tabular}{lll}
    \toprule
    Phase & Sign pattern & Classes \\
    \midrule
    Fig.~5: both axes IR-unstable, interior SCFT attracts
          & $(+,+)$      & $\{2\}$, and $\{1,9\}^{\ast}$ (nonSCFT) \\
    Mixed: one axis IR-unstable, one IR-stable
          & $(+,-)$ or $(-,+)$ & $\{4,16,21,23\}$ \\
    Fig.~6: both axes IR-stable, axis SCFTs attract
          & $(-,-)$      & \emph{none at $\Nf=0$} \\
    Marginal: $\alpha=0$ at both boundaries
          & $(0,0)$      & $\{44\}$ \\
    nonSCFT (no unitary IR at leading $N$)
          & ---          & $\{877\}$ and 20 flat-direction theories \\
    \bottomrule
  \end{tabular}
\end{center}

Thus, in the strict non-Veneziano limit, the classified
$\SU(N)\times\SU(N)$ landscape contains both the Figure~5 phase of
\cite{IW2005} and the mixed topology, but \emph{no} realisation of
the Figure~6 phase.  This matches the remark of \cite{IW2005} that
Figure~6 topologies were not found in the examples they studied.
Section~\ref{sec:phases} gives the phase portrait of each class in
detail and shows, as our main result, that a single-parameter
Veneziano deformation of class $16$ produces the first explicit
Figure~6 realisation in this landscape.

\subsection{Approximation: interior vs.\ boundary
  \texorpdfstring{$R_{\text{bif}}$}{R\_bif}}
\label{sec:boundary-caveat}

A technical point of prescription: the strict \cite{IW2005} recipe
computes $R_{\text{bif}}^{(a\star)}$ by $a$-maximising the
\emph{sub-quiver} on axis $a$ — node $a$ active with bifundamentals
promoted to $\Nf$-type flavours, node $b$ decoupled — rather than
using the interior fixed-point R-charge of the full two-node theory.
The two prescriptions in general differ because the sub-quiver is a
different CFT from the full two-node one.  In the tables of this
paper we use the interior $R_{\text{bif}}$ (from full-theory
$a$-maximisation) as a leading-$N$ proxy for the boundary
$R_{\text{bif}}^{(a\star)}$; both quantities share the sign of
$3R_{\text{bif}}-2$ under the mild assumption that neither
prescription has $R_{\text{bif}}$ cross the free-field value $2/3$,
and both yield the same qualitative phase assignment in
Table~\ref{tab:b0-boundary}.  The quantitative coefficient
$\alpha_{2}^{(1\star)}=-0.5730$ in class~16 is subject to
$\mathcal{O}(1/N)$ corrections under the strict boundary
prescription; this will be relevant for the Figure~6-window lower
edge discussed in Section~\ref{sec:phases}.
